\subsection{Interval Notation: Basics} 
\begin{definition}[Interval]
An \term{interval} is a set whose graph in the number line is a line segment or a ray.
%Some examples:
%\begin{icpic}
%\foreach \y in {0,-2,-4} {
	%\draw[axis] (-5,\y) -- (5,\y);
	%\draw[tick] (0,\y+0.2) -- (0,\y-0.2) node[anchor=north]{$0$};
	%\foreach \x in {1,2,...,4} {
		%\draw[tick] (\x,\y+0.2) -- (\x,\y-0.2) node[anchor=north]{$\x$};
		%\draw[tick] (-\x,\y+0.2) -- (-\x,\y-0.2) node[anchor=north]{$-\x\ $};
	%}
%}
%\draw[graphend] (-4,0.3) -- (1,0.3);
%\filldraw[dotempty] (-4,0.3) circle (0.1);
%\filldraw[dotfull] (1,0.3) circle (0.1);
%\draw[graphend,->] (-3,-1.7) -- (5,-1.7);
%\filldraw[dotfull] (-3,-1.7) circle (0.1);
%\draw[graphend,<-] (-5,-3.7) -- (2,-3.7);
%\filldraw[dotempty] (2,-3.7) circle (0.1);
%\end{icpic}
We can describe an interval in \term{interval notation}.
\end{definition}
\begin{displaybox}[Writing Interval Notation]
To describe an interval, we describe the left end, and seperately, describe the right end.
\begin{center}
\begin{tabular}{cc|c|cc}
\multicolumn{2}{c|}{Left End}&Comma&\multicolumn{2}{c}{Right End}\\\hline
Filled dot at \#&$\bigl[\#\bigr.$&
\multirow{3}{*}{\Large{,}}&
Filled dot at \#&$\bigl.\#\bigr]$\\
Empty dot at \#&$\bigl(\#\bigr.$&&
Empty dot at \#&$\bigl.\#)\bigr.$\\
Arrow (no end)&$\bigl(-\infty\bigr.$&&
Arrow (no end)&$\bigl.\infty\bigr)$
\end{tabular}
\end{center}
\end{displaybox}
\begin{Example} Write each set shown below in interval notation.
\begin{icpic}
\foreach \y in {0,-2,-4} {
	\numlinec[\y]{4}
	\draw (10,\y-0.2) node[anchor=base]{Interval Notation: \makeblank[1in]};
}
\graphoc[0]{-4}{1}
%\draw[graphend] (-4,0.3) -- (1,0.3);
%\filldraw[dotempty] (-4,0.3) circle (0.1);
%\filldraw[dotfull] (1,0.3) circle (0.1);
\graphca[-2]{-3}{5}
%\draw[graphend,->] (-3,-1.7) -- (5,-1.7);
%\filldraw[dotfull] (-3,-1.7) circle (0.1);
\graphao[-4]{-5}{2}
%\draw[graphend,<-] (-5,-3.7) -- (2,-3.7);
%\filldraw[dotempty] (2,-3.7) circle (0.1);
\end{icpic}
\end{Example}
\begin{Note}When writing interval notation, recall:
\begin{itemize}
\item We use the filled dot, and the square bracket, corresponding to the inequality $\leq$ or $\geq$; the number shown is \makeblank[1in] in the set.
\item We use the empty dot, and the round bracket, corresponding to the inequality $>$ or $<$; the number shown is \makeblank[1.5in] in the set.
\item The symbols $-\infty$ and $\infty$ are not \makeblank[1in], so the are \makeblank[1in] included in the set. They always have a \makeblank[1in] bracket.
\end{itemize}
\end{Note}
\begin{Example}
Sketch the graph of $-3<x\leq 5$, then describe the solution set in set-builder notation and in interval notation.
\begin{lpic}{}{4}
\regtextleft{1}{-3}{Set-Builder Notation:};
\regtextleft{1}{-4}{Interval Notation:};
\end{lpic}
\end{Example}
